\section{Data Sources and Acquisition}

Eight datasets were selected for this study, comprising four asset price series and four macroeconomic indicators. All data were obtained from publicly available sources. The sample period from January 1, 2020 to December 31, 2025 was chosen because it encompasses three distinct market regimes. The first was the COVID-19 crash and subsequent recovery from February to December 2020. The second was the inflation surge from 2021 through early 2023, during which US consumer price index year-on-year readings exceeded five percent. The third was the interest rate tightening cycle beginning in March 2022, during which the Federal Reserve raised the federal funds rate from near zero to over five percent. This tripartite regime structure provides a natural experiment for examining asset behaviour under contrasting macroeconomic conditions.

The US equity asset is Apple Incorporated (AAPL), listed on the NASDAQ exchange under ticker AAPL. As of 2025, Apple held a market capitalisation of approximately \$4.25 trillion, with a trailing price-to-earnings ratio of 35.08 and an estimated enterprise value of \$4.19 trillion \parencite{stockanalysis2025}. Apple was selected over alternative US technology stocks, particularly NVIDIA, because its valuation and business model present a fundamentally different risk profile that is instructive for a diversified portfolio analysis. NVIDIA derives over 80 percent of its revenue from data centre accelerator chips, exposing it heavily to cyclical capital expenditure cycles. Apple presented the opposite profile. Apple's approach to artificial intelligence has consistently prioritised on-device processing and user privacy over cloud-based deployment. The selection of AAPL over NVIDIA therefore captures a technology sector exposure that is diversified across hardware and services without the single-product concentration risk that characterises the pure-play semiconductor firms.

The Philippine equity asset is BDO Unibank Incorporated (BDO), the largest bank in the Philippines by total assets. As of December 2025, BDO reported total assets of PHP 5.27 trillion, maintaining its ranking as the top universal and commercial bank in the Philippines by asset size \parencite{bsp2025}. BDO was selected to provide exposure to the Philippine financial sector and to emerging market dynamics more broadly, which operate under different monetary policy cycles compared to developed markets.

The cryptocurrency asset is Bitcoin (BTC), the largest digital asset by market capitalisation. Bitcoin trades continuously on decentralised cryptocurrency exchanges and was included to evaluate the extent to which its return behaviour correlates with or diverges from traditional asset classes. The benchmark US equity index is represented by the SPDR S\&P 500 ETF Trust (SPY), which tracks the S\&P 500 index and is one of the most liquid exchange-traded funds globally.

Four macroeconomic series were obtained from the Federal Reserve Economic Data database maintained by the Federal Reserve Bank of St. Louis. The Consumer Price Index for All Urban Consumers (FRED series CPIAUCSL) measures headline consumer price inflation at monthly frequency. The 10-Year Treasury Constant Maturity Rate (FRED series DGS10) represents the yield on US government debt at the longest standard tenor and serves as a benchmark for risk-free rates. The Federal Funds Effective Rate (FRED series FEDFUNDS) captures the overnight rate at which depository institutions lend reserve balances to each other and reflects the stance of US monetary policy. The CBOE Volatility Index (FRED series VIXCLS) tracks the implied volatility of S\&P 500 index options over a 30-day forward period and is interpreted as a gauge of market fear and uncertainty.

Each dataset was downloaded from its respective source in CSV format and imported into R for processing. Asset price data for AAPL, BTC, and SPY were obtained from Yahoo Finance \parencite{yahoo_finance}. BDO data were sourced from Investing.com \parencite{investingcom}, as the Philippine stock data available through Yahoo Finance did not provide sufficient historical coverage for the required sample period. All four macroeconomic series were downloaded from the FRED database website \parencite{fred_data} using their native CSV export functionality. 

Please refer to the appendix for the complete R code demonstrating the data acquisition process.

Table~\ref{tab:data_acq} summarises the acquisition details for each dataset. The full sample period for each series extends from its earliest available observation through mid-2026. All series were subsequently restricted to the January 1, 2020 to December 31, 2025 analysis window during the data cleaning phase. No missing values or duplicate rows were detected in the raw downloads.

\begin{table}[H]
  \centering
  \caption{Data Acquisition Summary}
  \label{tab:data_acq}
  \begin{tabular}{l l r r r}
    \toprule
    Dataset & Source & {Original Observations} & {Missing Values} & {Duplicates} \\
    \midrule
    AAPL   & Yahoo Finance  & 2512 & 0 & 0 \\
    BDO    & Investing.com  & 2441 & 0 & 0 \\
    BTC    & Yahoo Finance  & 3652 & 0 & 0 \\
    SPY    & Yahoo Finance  & 2514 & 0 & 0 \\
    CPI    & FRED           & 952  & 0 & 0 \\
    DGS10  & FRED           & 16105 & 0 & 0 \\
    FEDFUNDS & FRED         & 863  & 0 & 0 \\
    VIX    & FRED           & 9216 & 0 & 0 \\
    \bottomrule
  \end{tabular}
\end{table}
